% =========================================================
%  ICDI Thai exam template, shared by every sitting folder
% =========================================================
% Each sitting folder (beforemidterm, midterm, prefinal, final) has a one-line
% groups.tex that \input's this file, so AMC still compiles <sitting>/groups.tex
% and every relative path below resolves from the sitting folder. What differs
% between sittings lives in <sitting>/config.tex: exam facts, \SelectTopic
% lines and the layout switches declared below. Edit this file once; never copy
% it into a sitting folder.
\documentclass[a4paper, twoside]{article}

% =========================================================
%  LAYOUT SWITCHES (set in config.tex with \<name>true)
% =========================================================
% Defaults give the midterm layout: TH Sarabun New, two-column questions, the
% back of the cover left blank, one choice per line.
\newif\ifCMUFont            % CMU font (../fonts/CMU) instead of TH Sarabun New
\newif\ifOneColumnQuestions % question pages in one column instead of two
\newif\ifShortChoicesOnLine % short choices on one line or in a 2x2 grid
\newif\ifQuestionsOnCoverBack % questions start on the back of the cover (saves a page per copy)

% =========================================================
%  📝 READ THE SITTING'S config.tex FIRST
% =========================================================
% config.tex sets the exam facts and switches, which the font and the packages
% below need, so it is read before anything else. Its \SelectTopic lines are
% only queued here and run after the question banks are loaded.
\newcommand{\ExamDraws}{}
\newcommand{\SelectTopic}[2]{%
    \expandafter\def\expandafter\ExamDraws\expandafter{\ExamDraws\DrawTopic{#1}{#2}}}
\input{config.tex}

%%% 1. Setup for Unicode and Thai Language %%%
\usepackage{fontspec}
\usepackage{xltxtra}
\usepackage{tikz}
\usepackage{multicol}
\usepackage{xcolor}
\usepackage{array}
\usepackage[most]{tcolorbox}

% --- FONT (before automultiplechoice: AMC sizes some spacing from it) ---
\ifCMUFont
% Light for body text and Regular for bold: CMU Regular as body text prints
% almost bold. License: ../fonts/CMU/CMU-FONT-licence-th.pdf.
\setmainfont[
    Path = ../fonts/CMU/,
    Extension = .ttf,
    UprightFont = *-Light,
    BoldFont = *-Regular,
    ItalicFont = *-Italic,
    BoldItalicFont = *-BoldItalic,
    Scale = 0.9
]{CMU}
\else
\setmainfont[
    Path = ../fonts/THSarabunNew/,
    Extension = .ttf,
    UprightFont = * ,
    BoldFont = * Bold,
    ItalicFont = * Italic,
    BoldItalicFont = * BoldItalic,
    Scale = 1.5
]{THSarabunNew}
\fi

% --- THAI SPACING ---
\XeTeXlinebreaklocale "th"
\XeTeXlinebreakskip = 0pt plus 1pt
\linespread{1.0}
\tolerance=5000
\emergencystretch=20pt

\usepackage[box,completemulti,separateanswersheet]{automultiplechoice}
% Package defaults (6 bits each) only encode values up to 64; this exam runs
% past that (page count, check code), so widen both to 9 bits (up to 511).
% @ must be a "letter" catcode to reach the package's internal macros.
\makeatletter
\def\AMC@NCBpage{9}
\def\AMC@NCBcheck{9}
\makeatother
\usepackage{totcount}

\usepackage{environ}   % for ../common/short-choices.tex
\ifShortChoicesOnLine% ---------------------------------------------------------------------------
% Short choices on one line (read by exam.tex when config.tex sets
% \ShortChoicesOnLinetrue). Every \begin{choices} is first run in a measuring
% pass; if all four choices fit on one line they are set with AMC's
% choiceshoriz, if two fit side by side they are set as a 2x2 grid, otherwise
% the usual one-per-line list is used. Needs the environ package (loaded by
% exam.tex).
% ---------------------------------------------------------------------------
\newlength{\ACmax}\newlength{\ACsum}\newlength{\ACtmp}\newcount\ACn
\newlength{\AClabel}\setlength{\AClabel}{2.6em}% per choice on one line: oval box and its gap, inside a .25\linewidth box
\newlength{\ACgridlabel}\setlength{\ACgridlabel}{2.6em}% per choice in the 2x2 grid
\newcommand{\ACmeasure}[1]{%
  \settowidth{\ACtmp}{#1}%
  \global\advance\ACn by 1
  \global\advance\ACsum by \ACtmp
  \ifdim\ACtmp>\ACmax \global\ACmax=\ACtmp\fi}
\let\ACorigchoices=\choices
\let\endACorigchoices=\endchoices
\newcommand{\ACboxw}{.5\linewidth}% box width per choice: .25 on one line, .5 in the grid
\def\AMCanswer#1#2{\makebox[\ACboxw][l]{#1\hspace*{\AMChorizBoxSep}#2}\allowbreak}
\RenewEnviron{choices}[1][r]{%
  \global\ACmax=0pt \global\ACsum=0pt \global\ACn=0
  \begingroup
    \let\correctchoice=\ACmeasure
    \let\wrongchoice=\ACmeasure
    \setbox0=\vbox{\BODY}%
  \endgroup
  \setlength{\ACtmp}{\ACmax}\addtolength{\ACtmp}{\AClabel}%
  \ifdim\ACtmp<.25\linewidth
    \renewcommand{\ACboxw}{.25\linewidth}%
    \par\vspace{.6ex}\noindent\begin{choicescustom}[#1]\BODY\end{choicescustom}\par
  \else
    \setlength{\ACtmp}{\ACmax}\addtolength{\ACtmp}{\ACgridlabel}%
    \ifdim\ACtmp<.5\linewidth
      \renewcommand{\ACboxw}{.5\linewidth}%
      \par\vspace{.6ex}\noindent\begin{choicescustom}[#1]\BODY\end{choicescustom}\par
    \else
      \begin{ACorigchoices}[#1]\BODY\end{ACorigchoices}%
    \fi
  \fi}
\fi
\regtotcounter{AMCquestionaff}

% =========================================================
%  🔧 TOPIC SELECTION (DO NOT EDIT)
% =========================================================
% \SelectTopic{group}{count} → copies N random questions
% from the given group into the "selected" pool. The group is shuffled once
% here, outside \onecopy, so every copy gets the SAME N questions (a random
% pick, not the first N in file order); \shufflegroup{selected} inside
% \onecopy then gives each copy its own order.
\newcommand{\DrawTopic}[2]{%
    \shufflegroup{#1}%
    \copygroup[#2]{#1}{selected}%
}

% =========================================================
%  📚 LOAD ALL QUESTION BANKS, THEN DRAW config.tex's \SelectTopic LINES
% =========================================================
% Shared bank lives one level up, in questions/ — every sitting
% folder draws from it.
% One line per chapter file. Add a line when you add a file to questions/.
% ======================================================================
% Chapter 1: Digital economy basics (SAMPLE questions for the guide)
% Group: ch01_basics      ID prefix: DE
% One question per \element. Every \element in this file uses the same group.
% ======================================================================

\element{ch01_basics}{
  \begin{question}{DE001}
    ร้านกาแฟเปลี่ยนจากบัตรสะสมแต้มกระดาษมาเป็นแต้มในแอปของร้าน การเปลี่ยนแปลงนี้คือข้อใด
    \begin{choices}
      \correctchoice{Digitization}
      \wrongchoice{Outsourcing}
      \wrongchoice{Franchising}
      \wrongchoice{Crowdfunding}
    \end{choices}
    \explain{\textbf{คำอธิบาย:} การเปลี่ยนข้อมูลที่อยู่บนกระดาษให้อยู่ในรูปดิจิทัลคือ Digitization}
  \end{question}
}

\element{ch01_basics}{
  \begin{question}{DE002}
    แอปเรียกรถจับคู่ผู้โดยสารกับคนขับโดยไม่มีรถเป็นของตนเอง ธุรกิจนี้เป็นแบบใด
    \begin{choices}
      \wrongchoice{ผู้ผลิตสินค้า}
      \correctchoice{แพลตฟอร์ม}
      \wrongchoice{ผู้ค้าส่ง}
      \wrongchoice{ผู้ให้เช่ารถ}
    \end{choices}
    \explain{\textbf{คำอธิบาย:} แพลตฟอร์มสร้างมูลค่าจากการจับคู่สองฝั่งตลาด ไม่ได้เป็นเจ้าของสินทรัพย์หลักเอง}
  \end{question}
}

\element{ch01_basics}{
  \begin{question}{DE003}
    ยิ่งมีผู้ใช้แอปแชตมากขึ้น แอปก็ยิ่งมีประโยชน์กับผู้ใช้แต่ละคนมากขึ้น ปรากฏการณ์นี้คือข้อใด
    \begin{choices}
      \wrongchoice{Economies of scope}
      \wrongchoice{Price discrimination}
      \correctchoice{Network effect}
      \wrongchoice{Moral hazard}
    \end{choices}
  \end{question}
}

\element{ch01_basics}{
  \begin{question}{DE004}
    เจ้าของคอนโดนำห้องว่างไปให้นักท่องเที่ยวเช่ารายวันผ่านแอป ข้อใดคือทรัพยากรที่ถูกนำมาใช้ให้เกิดประโยชน์
    \begin{choices}
      \correctchoice{สินทรัพย์ที่ว่างอยู่}
      \wrongchoice{แรงงานที่มีทักษะสูง}
      \wrongchoice{เงินทุนจากธนาคาร}
      \wrongchoice{สิทธิบัตรของบริษัท}
    \end{choices}
  \end{question}
}

\element{ch01_basics}{
  \begin{question}{DE005}
    บริษัทขายเพลงแบบจ่ายรายเดือนแทนการขายเป็นอัลบั้ม รูปแบบรายได้นี้คือข้อใด
    \begin{choices}
      \wrongchoice{Advertising}
      \correctchoice{Subscription}
      \wrongchoice{Commission}
      \wrongchoice{Licensing}
    \end{choices}
  \end{question}
}

\element{ch01_basics}{
  \begin{question}{DE006}
    ร้านค้าออนไลน์ขายสินค้าให้ผู้บริโภคทั่วไปโดยตรง เป็นธุรกรรมประเภทใด
    \begin{choices}
      \wrongchoice{B2B}
      \correctchoice{B2C}
      \wrongchoice{C2C}
      \wrongchoice{G2B}
    \end{choices}
  \end{question}
}

\element{ch01_basics}{
  \begin{question}{DE007}
    ต้นทุนของการทำสำเนาไฟล์ e-book เพิ่มอีกหนึ่งชุดใกล้เคียงข้อใดมากที่สุด
    \begin{choices}
      \correctchoice{เกือบเป็นศูนย์}
      \wrongchoice{เท่ากับต้นทุนพิมพ์หนังสือหนึ่งเล่ม}
      \wrongchoice{สูงขึ้นตามจำนวนสำเนา}
      \wrongchoice{เท่ากับค่าลิขสิทธิ์ทั้งหมด}
    \end{choices}
  \end{question}
}

\element{ch01_basics}{
  \begin{question}{DE008}
    ผู้ขายของมือสองโพสต์ขายให้ผู้ซื้อที่เป็นบุคคลทั่วไปบนเว็บประกาศ เป็นธุรกรรมประเภทใด
    \begin{choices}
      \wrongchoice{B2B}
      \wrongchoice{B2C}
      \correctchoice{C2C}
      \wrongchoice{B2G}
    \end{choices}
  \end{question}
}

\element{ch01_basics}{
  \begin{question}{DE009}
    แอปแผนที่ให้ใช้ฟรีแต่มีรายได้จากร้านค้าที่จ่ายเงินให้แสดงตำแหน่งเด่น รูปแบบรายได้นี้คือข้อใด
    \begin{choices}
      \correctchoice{Advertising}
      \wrongchoice{Subscription}
      \wrongchoice{Pay per use}
      \wrongchoice{Licensing}
    \end{choices}
  \end{question}
}

\element{ch01_basics}{
  \begin{question}{DE010}
    ธนาคารปิดสาขาบางแห่งเพราะลูกค้าส่วนใหญ่ทำธุรกรรมผ่านแอปแทน ข้อใดอธิบายเหตุการณ์นี้ได้ดีที่สุด
    \begin{choices}
      \wrongchoice{Digitization}
      \wrongchoice{Crowdsourcing}
      \correctchoice{Digital transformation}
      \wrongchoice{Vertical integration}
    \end{choices}
  \end{question}
}

% ======================================================================
% Chapter 2: Data and cloud (SAMPLE questions for the guide)
% Group: ch02_cloud      ID prefix: DC
% ======================================================================

\element{ch02_cloud}{
  \begin{question}{DC001}
    บริษัทเช่าเครื่องเสมือนจากผู้ให้บริการคลาวด์ แล้วติดตั้งระบบปฏิบัติการและโปรแกรมเอง เป็นบริการแบบใด
    \begin{choices}
      \correctchoice{IaaS}
      \wrongchoice{PaaS}
      \wrongchoice{SaaS}
      \wrongchoice{On-premise}
    \end{choices}
  \end{question}
}

\element{ch02_cloud}{
  \begin{question}{DC002}
    สำนักงานใช้อีเมลและเอกสารออนไลน์ผ่านเว็บเบราว์เซอร์ โดยไม่ต้องติดตั้งหรือดูแลเซิร์ฟเวอร์ เป็นบริการแบบใด
    \begin{choices}
      \wrongchoice{IaaS}
      \wrongchoice{PaaS}
      \correctchoice{SaaS}
      \wrongchoice{Colocation}
    \end{choices}
  \end{question}
}

\element{ch02_cloud}{
  \begin{question}{DC003}
    ทีมพัฒนาส่งเฉพาะโค้ดขึ้นระบบ ส่วนผู้ให้บริการดูแลเซิร์ฟเวอร์และระบบปฏิบัติการให้ เป็นบริการแบบใด
    \begin{choices}
      \wrongchoice{IaaS}
      \correctchoice{PaaS}
      \wrongchoice{SaaS}
      \wrongchoice{On-premise}
    \end{choices}
  \end{question}
}

\element{ch02_cloud}{
  \begin{question}{DC004}
    ข้อใดเป็นข้อมูลแบบไม่มีโครงสร้าง (unstructured data)
    \begin{choices}
      \wrongchoice{ตารางยอดขายรายวัน}
      \wrongchoice{รายชื่อนักศึกษาพร้อมรหัส}
      \correctchoice{คลิปเสียงสายลูกค้า}
      \wrongchoice{ตารางราคาสินค้า}
    \end{choices}
  \end{question}
}

\element{ch02_cloud}{
  \begin{question}{DC005}
    ร้านค้าพบว่าที่อยู่ลูกค้าจำนวนมากสะกดผิดและซ้ำกัน เป็นปัญหาด้านใดของข้อมูล
    \begin{choices}
      \wrongchoice{Volume}
      \wrongchoice{Velocity}
      \wrongchoice{Variety}
      \correctchoice{Veracity}
    \end{choices}
  \end{question}
}

\element{ch02_cloud}{
  \begin{question}{DC006}
    เว็บขายตั๋วคอนเสิร์ตต้องรับผู้ใช้มากขึ้นสิบเท่าในวันเปิดขาย คุณสมบัติใดของคลาวด์ช่วยได้มากที่สุด
    \begin{choices}
      \correctchoice{ขยายทรัพยากรได้ตามความต้องการ}
      \wrongchoice{ข้อมูลอยู่ในประเทศเดียวกับผู้ใช้}
      \wrongchoice{ค่าบริการคงที่ทุกเดือน}
      \wrongchoice{ไม่ต้องใช้อินเทอร์เน็ต}
    \end{choices}
  \end{question}
}

\element{ch02_cloud}{
  \begin{question}{DC007}
    เซ็นเซอร์ในโรงงานส่งค่าอุณหภูมิเข้ามาทุกวินาที เป็นลักษณะใดของข้อมูล
    \begin{choices}
      \wrongchoice{Volume}
      \correctchoice{Velocity}
      \wrongchoice{Variety}
      \wrongchoice{Value}
    \end{choices}
  \end{question}
}

\element{ch02_cloud}{
  \begin{question}{DC008}
    ตัวเลข 36.5 ที่ไม่มีหน่วยและไม่รู้ว่าวัดอะไร อยู่ในระดับใดของ DIKW
    \begin{choices}
      \correctchoice{Data}
      \wrongchoice{Information}
      \wrongchoice{Knowledge}
      \wrongchoice{Wisdom}
    \end{choices}
  \end{question}
}

\element{ch02_cloud}{
  \begin{question}{DC009}
    บริษัทจ่ายค่าคลาวด์ตามชั่วโมงที่เปิดเครื่องใช้งานจริง เป็นรูปแบบการคิดค่าบริการแบบใด
    \begin{choices}
      \wrongchoice{เหมาจ่ายรายปี}
      \correctchoice{จ่ายตามการใช้งาน}
      \wrongchoice{ซื้อขาดครั้งเดียว}
      \wrongchoice{ค่าธรรมเนียมแรกเข้า}
    \end{choices}
  \end{question}
}

\element{ch02_cloud}{
  \begin{question}{DC010}
    โรงพยาบาลเก็บข้อมูลผู้ป่วยไว้บนเซิร์ฟเวอร์ของตนเอง และใช้คลาวด์สาธารณะสำหรับเว็บไซต์ เป็นรูปแบบใด
    \begin{choices}
      \wrongchoice{Public cloud}
      \wrongchoice{Private cloud}
      \correctchoice{Hybrid cloud}
      \wrongchoice{Community cloud}
    \end{choices}
  \end{question}
}

\ExamDraws

% --- AMC FORMATTING ---
\def\AMCbeginQuestion#1#2{\par\noindent\textbf{#1.}~#2}
\def\AMCformQuestion#1{\vspace{-0.15cm}\par\noindent\makebox[1.9em][r]{\textbf{#1.}}}
\setdefaultgroupmode{withoutreplacement}

% --- AMC BOX STYLE ---
% Absolute sizes (the TH Sarabun 1.5x values of 2.5ex/.4ex/.1em/.5em/.3em), so the
% bubbles and the ID grid keep the same geometry whatever main font is loaded;
% \IDWriteBoxes below depends on the 7.29 mm column pitch this gives.
\AMCboxStyle{shape=oval,size=12.75pt,down=2.04pt,rule=.5pt,outsidesep=1.5pt,color=black}
\AMCcodeHspace=7.5pt
\AMCcodeVspace=7.5pt
\AMCformHSpace=4.5pt

% --- COVER PAGE SHADES (black and white only; the exam is printed in mono) ---
\definecolor{coverDark}{gray}{0}
\definecolor{coverGray}{gray}{0.3}
\definecolor{coverRule}{gray}{0.55}
\definecolor{coverLight}{gray}{1}   % no background fills: saves toner
\newcommand{\CoverCell}[3]{%
    \begin{tabular}{@{}c@{}}{\small\color{coverGray} #1}\\[-0.6ex]{\footnotesize\color{coverGray} #2}\\[0.4ex]{\large\bfseries #3}\end{tabular}}
\newcommand{\AnsCell}[3]{%
    \begin{tabular}{@{}c@{}}{\footnotesize\color{coverGray} #1}\\[-0.8ex]{\scriptsize\color{coverGray} #2}\\[0.2ex]{\normalsize\bfseries #3}\end{tabular}}
% English title and date, set per sitting in config.tex.
\providecommand{\ExamTitleEN}{}
\providecommand{\ExamDateEN}{}
\providecommand{\ExamDateNum}{\ExamDate}
% Thai date with the English date under it (cover and answer-sheet cells).
\newcommand{\ExamDateBoth}[1]{%
    \begin{tabular}{@{}c@{}}\ExamDate\\[-0.3ex]{#1\mdseries\color{coverGray}\ExamDateEN}\end{tabular}}
\providecommand{\CourseName}{Course Name}   % override in config.tex if needed

% --- QUESTION PAGES ---
\setlength{\AMCinterBquest}{3.2ex}
\setlength{\columnseprule}{0.4pt}
\renewcommand{\columnseprulecolor}{\color{coverRule}}

% Footer on the question pages only (not the cover, not the answer sheet), centered: course | exam | page.
% AMC's page style owns the header/footer slots (corner marks), so the footer is
% drawn in the shipout foreground layer instead and never touches the marks.
\newif\ifExamFooter
\AddToHook{shipout/foreground}{%
  \ifExamFooter\ifnum\value{page}>1
    \put(\dimexpr 1in+\oddsidemargin\relax,-\dimexpr\paperheight-1.7cm\relax){%
      \makebox[\textwidth][c]{\small\color{coverGray}%
        \CourseCode\ \CourseName\ \ \textbar\ \ \ExamTitle\ \ \textbar\ \ หน้า \thepage}}%
  \fi\fi}

% Nine write-in boxes, one above each column of the \AMCcode{id}{9} grid.
% The grid's column pitch is 7.29 mm and each circle is 4.66 mm wide (measured
% from data/layout.sqlite), so each box is centered 1.57 mm into its 7.29 mm cell (checked on a 200 dpi render);
% centering this picture and the grid in the same box lines them up.
\newcommand{\IDWriteBoxes}{%
  \begin{tikzpicture}
    \useasboundingbox (0,0) rectangle (9*7.29mm, 7.5mm);
    \foreach \i in {0,...,8}
      \draw[coverDark, line width=0.5pt, rounded corners=0.6mm]
        (\i*7.29mm+1.57mm-2.9mm, 0) rectangle ++(5.8mm, 7.5mm);
  \end{tikzpicture}}

% Boxed section title used on the answer sheet (same look as the cover's boxes).
\tcbset{examtitled/.style={enhanced, colback=white, colframe=coverDark, boxrule=0.5pt, arc=2mm,
    left=3mm, right=3mm, top=3mm, bottom=2mm, fonttitle=\small\bfseries,
    coltitle=coverDark, colbacktitle=white,
    attach boxed title to top left={xshift=4mm, yshift=-2.2mm},
    boxed title style={arc=1.5mm, boxrule=0.5pt, colframe=coverDark}}}


% Framed answer-sheet boxes (name, student ID). AMC records every bubble and
% zone position with \pdfsavepos, which cannot see the coordinate shift
% tcolorbox's enhanced skin applies to its content (it draws it inside a TikZ
% node), so an \AMCcode or \namefield placed inside the box is detected ~4 mm
% off in the scan. The content is therefore typeset into a saved box *outside*
% the tcolorbox, and the empty frame is drawn first, behind it.
% #1 title, #2 extra space below the content, #3 content.
\newsavebox{\SheetBoxContent}
\newcommand{\SheetBox}[3]{%
  \sbox{\SheetBoxContent}{%
    \begin{minipage}[t]{\dimexpr\linewidth-6mm-1pt\relax}#3\end{minipage}}%
  \par\noindent
  \begin{tcolorbox}[examtitled, title={#1},
      height=\dimexpr\ht\SheetBoxContent+\dp\SheetBoxContent+5mm+#2+1pt\relax]
  \end{tcolorbox}%
  \par\nointerlineskip
  \vspace{-\dimexpr\ht\SheetBoxContent+\dp\SheetBoxContent+2mm+#2+0.5pt\relax}%
  \noindent\hspace*{\dimexpr3mm+0.5pt\relax}\usebox{\SheetBoxContent}\par
  \vspace{\dimexpr2mm+#2\relax}}

\newcommand{\IDSheetBox}{%
  \SheetBox{รหัสนักศึกษา / Student ID}{3mm}{%
    \centering\vspace{0.5mm}
    {\footnotesize เขียนรหัสลงในช่อง แล้วฝนวงกลมใต้ช่องให้ตรงกัน}\\[-0.4ex]
    {\footnotesize\color{coverGray} Write your ID, then fill the circles below.}\\[1mm]
    \IDWriteBoxes\\[1.2mm]
    \AMCcode{id}{9}}}

\newcommand{\NameSheetBox}{%
  \SheetBox{ชื่อ - นามสกุล / Name}{2mm}{%
    \namefield{\begin{minipage}{\linewidth}
        \vspace{5mm}
        \dotfill\par
        \vspace{1mm}
    \end{minipage}}}}

\begin{document}
\raggedbottom

% Start of AMC exam generation
\onecopy{\NumCopies}{

    % =========================================================
    % 1. EXAM COVER PAGE
    % =========================================================
    \shufflegroup{selected}
    \global\ExamFootertrue

    % --- Header: college logo and name ---
    \noindent
    \begin{minipage}[c]{2.3cm}
        \includegraphics[width=2.3cm]{../logos/logo-ICDI-gray.png}
    \end{minipage}\hspace{0.4cm}%
    \begin{minipage}[c]{\dimexpr\linewidth-2.7cm\relax}
        {\Large\bfseries วิทยาลัยนานาชาตินวัตกรรมดิจิทัล มหาวิทยาลัยเชียงใหม่}\\[-0.2ex]
        {\color{coverGray}International College of Digital Innovation, Chiang Mai University}
    \end{minipage}

    \vspace{0.5cm}

    % --- Title band ---
    \begin{tcolorbox}[enhanced, colback=white, colframe=coverDark, arc=2mm,
        boxrule=0.8pt, left=6mm, right=6mm, top=2.5mm, bottom=2.5mm, halign=center]
        {\XeTeXlinebreakskip=0pt\relax\fontsize{24}{28}\selectfont\bfseries \ExamTitle}\\[1.2ex]
        {\Large\color{coverGray} \ExamTitleEN}\\[0.6ex]
        {\large \CourseCode\ \ \CourseName}
    \end{tcolorbox}

    \vspace{0.3cm}

    % --- Exam facts: five cells ---
    \begin{tcolorbox}[enhanced, colback=coverLight, colframe=coverRule, boxrule=0.5pt,
        arc=2mm, left=3mm, right=3mm, top=1.2mm, bottom=1.6mm]
        \begin{tabular*}{\linewidth}{@{\extracolsep{\fill}}ccccc@{}}
            \CoverCell{กระบวนวิชา}{Course}{\CourseCode} &
            \CoverCell{ตอน}{Section}{\CourseSection} &
            \CoverCell{วันที่สอบ}{Date}{\ExamDateBoth{\small}} &
            \CoverCell{เวลา}{Time}{\ExamTime} &
            \CoverCell{จำนวนข้อสอบ}{No.\ of questions}{\total{AMCquestionaff}}
        \end{tabular*}
    \end{tcolorbox}

    \vspace{0.3cm}

    % --- Student details ---
    \begin{tcolorbox}[enhanced, colback=white, colframe=coverDark, boxrule=0.5pt, arc=2mm,
        left=4mm, right=4mm, top=3mm, bottom=3mm,
        title={\bfseries ข้อมูลนักศึกษา / Student information}, fonttitle=\normalsize,
        coltitle=coverDark, colbacktitle=white, attach boxed title to top left={xshift=4mm, yshift=-2.2mm},
        boxed title style={arc=1.5mm, boxrule=0.5pt, colframe=coverDark}]
        \vspace{1.5mm}
        ชื่อ - นามสกุล / Name\enspace\dotfill\par
        \vspace{3mm}
        \noindent รหัสนักศึกษา / Student ID\enspace
        \begin{tikzpicture}[baseline=1.5mm]
            \foreach \i in {0,...,8} {\draw[coverDark, line width=0.5pt] (\i*0.62,0) rectangle ++(0.56,0.62);}
        \end{tikzpicture}%

    \end{tcolorbox}

    \vspace{0.2cm}

    % --- Instructions ---
    \begin{tcolorbox}[enhanced, colback=white, colframe=coverRule, boxrule=0.5pt, arc=2mm,
        left=3mm, right=4mm, top=3mm, bottom=2mm,
        title={\bfseries คำชี้แจง / Instructions}, fonttitle=\normalsize,
        coltitle=coverDark, colbacktitle=white, attach boxed title to top left={xshift=4mm, yshift=-2.2mm},
        boxed title style={arc=1.5mm, boxrule=0.5pt, colframe=coverDark}]
        \vspace{1mm}
        \linespread{1.3}\selectfont
        \begin{enumerate}\setlength{\itemsep}{3pt}\setlength{\parskip}{0pt}
            \item ข้อสอบมีจำนวนทั้งสิ้น \textbf{\total{AMCquestionaff} ข้อ} ให้ฝนคำตอบลงใน\textbf{กระดาษคำตอบ}ที่แยกไว้ท้ายเล่มเท่านั้น
            \item ก่อนเข้าห้องสอบให้กรรมการคุมสอบประกาศให้นักศึกษาทราบว่า สิ่งของใดที่อนุญาตให้นำเข้าห้องสอบ และสิ่งของใดที่ห้ามนำเข้าห้องสอบ
            \item ให้นักศึกษาเขียนชื่อ - นามสกุล และรหัสประจำตัว ให้ชัดเจน
            \item ก่อนเริ่มดำเนินการสอบให้กรรมการคุมสอบประกาศในห้องสอบให้นักศึกษาทราบเกี่ยวกับชื่อกระบวนวิชาที่จัดสอบ
            \item ให้นักศึกษาปฏิบัติตามคำชี้แจงของกรรมการคุมสอบอย่างเคร่งครัด
            \item ให้ปฏิบัติตามข้อบังคับมหาวิทยาลัยเชียงใหม่ ว่าด้วยการสอบของนักศึกษา พ.ศ. 2554 อย่างเคร่งครัด
        \end{enumerate}
    \end{tcolorbox}

    \vspace{0.15cm}

    % --- Integrity warning ---
    \begin{tcolorbox}[enhanced, colback=coverLight, colframe=coverDark, boxrule=0pt,
        leftrule=0.6mm, sharp corners, left=4mm, right=4mm, top=1.5mm, bottom=1.5mm]
        \linespread{1.3}\selectfont
        \textbf{คำเตือน:} การทุจริตหรือกระทำการใด ๆ ซึ่งเป็นการทุจริตในการสอบ
        \textbf{เป็นความผิดวินัยร้ายแรง} และ\textbf{จะได้รับการพิจารณาให้อักษรลำดับขั้น F หรือ U}
    \end{tcolorbox}

    \vspace{0.4cm}

    % --- Agreement ---
    \begin{center}
        \textbf{การเข้าสอบถือว่านักศึกษายอมรับและจะปฏิบัติตามคำชี้แจงข้างต้นทุกประการ}\\[0.2ex]
        {\color{coverGray}By taking this exam, you agree to the terms above.}
    \end{center}

    \vfill

    \ifQuestionsOnCoverBack
        \clearpage
    \else
        \AMCcleardoublepage
    \fi

    % =========================================================
    % 2. QUESTIONS SECTION
    % =========================================================

    \ifOneColumnQuestions
        \insertgroup{selected}
    \else
        \setlength{\columnsep}{.8cm}
        \begin{multicols}{2}
            \insertgroup{selected}
        \end{multicols}
    \fi

    \cleardoublepage

    % =========================================================
    % 3. ANSWER SHEET SECTION
    % =========================================================
    \AMCformBegin
    \global\ExamFooterfalse

    % --- Title band: sheet name, then the exam on the next line ---
    \begin{tcolorbox}[enhanced, colback=white, colframe=coverDark, arc=2mm, boxrule=0.8pt,
        left=5mm, right=5mm, top=1mm, bottom=1.2mm]
        {\XeTeXlinebreakskip=0pt\relax\fontsize{16}{19}\selectfont\bfseries กระดาษคำตอบ}
        \hspace{0.5em}{\large\color{coverGray} ANSWER SHEET}\\[-0.2ex]
        {\normalsize\bfseries \ExamTitle}\hspace{0.5em}{\normalsize\color{coverGray}\ExamTitleEN}
    \end{tcolorbox}

    \vspace{-0.1cm}
    \begin{tcolorbox}[enhanced, colback=coverLight, colframe=coverRule, boxrule=0.5pt,
        arc=2mm, left=3mm, right=3mm, top=0.8mm, bottom=1.2mm]
        \begin{tabular*}{\linewidth}{@{\extracolsep{\fill}}ccccc@{}}
            \AnsCell{กระบวนวิชา}{Course}{\CourseCode} &
            \AnsCell{ตอน}{Section}{\CourseSection} &
            \AnsCell{วันที่สอบ}{Date}{\ExamDateNum} &
            \AnsCell{เวลา}{Time}{\ExamTime} &
            \AnsCell{จำนวนข้อสอบ}{No.\ of questions}{\total{AMCquestionaff}}
        \end{tabular*}
    \end{tcolorbox}

    \vspace{0.15cm}

    % --- Left: name, then instructions. Right: student ID grid. ---
    \noindent
    \begin{minipage}[t]{0.47\linewidth}
        \vspace{0pt}
        \NameSheetBox

        \vspace{0.1cm}
        \begin{tcolorbox}[examtitled, title={คำชี้แจง / Instructions}]
            \vspace{0.5mm}\raggedright\small
            \textbf{ดึงกระดาษคำตอบแผ่นนี้ออกจากเล่ม}\\
            {\color{coverGray}Pull this answer sheet (the last page) out of the booklet.}

            \vspace{1.5mm}
            ใช้ปากกาหรือดินสอ 2B ฝนวงกลมให้\textbf{เต็มวง}\\
            {\color{coverGray}Use a pen or a 2B pencil. Fill the circle completely.}

            \vspace{2mm}
            \centering
            \begin{tabular}{c@{\hspace{1.2cm}}c}
                \begin{tikzpicture}[baseline=-0.5ex]
                    \fill (0,0) circle (0.2);
                \end{tikzpicture} &
                \begin{tikzpicture}[baseline=-0.5ex]
                    \draw (0,0) circle (0.2);
                    \draw[line width=0.6pt] (-0.14,0.14) -- (0.14,-0.14);
                    \draw[line width=0.6pt] (-0.14,-0.14) -- (0.14,0.14);
                \end{tikzpicture} \\[0.8ex]
                \textbf{ใช้ได้} {\small\color{coverGray}Valid} &
                \textbf{ใช้ไม่ได้} {\small\color{coverGray}Invalid}
            \end{tabular}
        \end{tcolorbox}
    \end{minipage}\hfill
    \begin{minipage}[t]{0.5\linewidth}
        \vspace{0pt}
        \IDSheetBox
    \end{minipage}

    \vspace{0.25cm}
    {\color{coverDark}\hrule height 0.8pt}
    \vspace{0.1cm}

    \setlength{\columnsep}{0.5cm}
    \begin{multicols}{3}
        \AMCform
    \end{multicols}

    % Pad every copy to an even number of pages, or duplex printing puts the
    % next copy on the back of this one's last page (the 180-question final
    % runs to 29 or 31 pages). The question part already starts the answer
    % sheet on a front side with \cleardoublepage above.
    \clearpage
    \ifodd\value{page}\else\null\clearpage\fi
} 

\end{document}
